\scanpage{216}
\mspage{450}
% archive sheet number: 215
(Ce n'est pas \inserted{\handwritten{certes}} la première fois \inserted{\handwritten{ni la dernière que je suis ``soufflé''}} dans cette brillante opération \char34{}au service de la science\char34{}\handwritten{...)} \deleted{\unsure{même|ni même} \illegible{3}} J'apprends par sa lettre (datée du 15 février 85) que conformément \deleted{\unsure{aux|au}} \char34{}à la façon habituelle d'agir quand un ouvrage contient des contributions de plusieurs auteurs\char34{} (sic), il n'y avait pas lieu de s'adresser plus spécialement à moi, qui n'étais que \uline{directeur} du Séminaire ... \deleted{\unsure{(}}Les cinq \char34{}auteurs\char34{} de SGA 5 sont Bucur, Houzel, Illusie, Jouanolou, Serre, à l'exclusion de ma modeste personne, qui \deleted{\unsure{ne}} ne figure que comme \char34{}directeur\char34{} - sans doute purement honorifique, en \inserted{\handwritten{av}}ai\inserted{\handwritten{s}}-je trop dit (\textsuperscript{\unsure{x|*}}) - pour ce brillant séminaire.

Dès avant de recevoir cette lettre instructive et trouvant le temps long (n'ayant rien reçu pendant un mois), j'ai pris ma plus belle plume (en allemand) pour écrire au Dr. K.F Springer en personne, \inserted{\handwritten{lequel}} fai\deleted{sant}\handwritten{t} partie des directeurs responsables de la maison Springer. Ça a été une belle \deleted{let} lettre de deux pages machine, lui expliquant que j'étais bien peiné d'une \unsure{L|l}ongue série de désagréments dans ma relation à la maison Springer, et au delà de celles-ci, d'un certain nombre d\deleted{\unsure{1}}irrégularités grossières à mon égard, dont je me contentais pour le moment de lui soumettre deux, qui me p\deleted{\unsure{e}}araissaient particulièrement flagrant\deleted{\unsure{s}}\handwritten{es}: la publication de deux volumes des Lecture Notes (n\textsuperscript{\uline{o}}s 288,589) parus sous mon nom et sans \deleted{\unsure{aucun}} juger nécessaire de me consulter. Que dans ces deux \deleted{cas}\inserted{textes}, les idées, méthodes et résultats que j'avais développés dans les séminaires ora\uline{u}x, étaient raccourcis ou mutilés parfois au point d'être méconnaissables. Que la coïncidence de ce dernier fait, \handwritten{avec} \deleted{d}\handwritten{l}es circonstances inhabituelles qui ont entouré la publication de ces deux volumes, ne pouvait être pour moi l'effet d'un pur hasar\deleted{\unsure{s}}d. Et que je m'attendais à des excuses publiques\deleted{\unsure{d} de la part} et sans réserve de la part de la maison\handwritten{/}Springer, sous une forme qui serait à fixer d'un commun accord, une fois \deleted{que la question}\inserted{qu'un accord} de principe \deleted{ét}\handwritten{ser}ait acquis. Que j'espérais qu'il aurait à coeur comme moi de mettre fin à une situation déplaisante et inadmissible et de trouver une solution\unsure{n}qui soit à la hauteur des circonstances (\char34{}eine dem Fall geziemende Lösung zu finden\char34{}, ce qui \deleted{est} est même plus distingué encore), \char34{}ho\deleted{\unsure{x}}chachtungsvoll\char34{} (comme il se doit) signé de ma plus belle main.

Pour jouer cartes sur table, il me semble que j'ai joué cartes sur table ! Là il ne pourra pas dire, Monsieur K.F. S\handwritten{p}ringer, \deleted{\unsure{c}}\handwritten{q}u'il \deleted{\unsure{n'a}}\handwritten{n'a} pas été informé personellement de la situation, et de première main enco\deleted{\unsure{re}}re\handwritten{,}\deleted{\unsure{é}} par nul autre que le principal intéressé lui-mê\handwritten{m}e!

Le hasard fait bien les choses: j'ai \deleted{réçu} fini par recevoir une réponse (un bon mois après\inserted{\handwritten{ encore}}) pas\unsure{,|;}plus tard que hier. C'est si court que je ne
